%% kapa: exam=1 round=37 year=2026 period=1 form=A subject=economics subject_ko=경제학원론 seq=28 number=68
%% source: Q-Net 감정평가사 기출문제 (cst00309) artlSeq=5258924 file=raw/1008/5258924/2026년 제37회 감정평가사 1차 1교시.pdf
%% conversion: kordoc 4.24.51 -> tools/gen1.py (auto); review=proofread-sonnet; flags=none
\begin{question}{68}
유동성함정에 관한 설명으로 옳지 않은 것은?
\begin{choices}
  \item 화폐수요의 이자율 탄력성이 무한대(∞)이다.
  \item 정부지출 증가로 인한 구축효과가 나타난다.
  \item 조세감면을 통한 재정정책이 소득증가에 효과적이다.
  \item 통화공급이 증가하더라도 이자율이 더 이상 낮아지지 않는다.
  \item 확장적 통화정책을 시행하더라도 소득에 변화가 없다.
\end{choices}
\end{question}
\answer{2}
